\scanpage{319}

\mspage{58}
% archive sheet number: 318

de $f_i$ est l'élément $1$ de $\mathrm{Ty}$ (correspondant à l'intervalle ponctuel), et comme $\operatorname{Hom}(1,\tau)=\varnothing$ pour $\tau\neq1$, cela signifie que les flèches qui suivent $f_i$ doivent être toutes égales à la flèche $\mathrm{id}_1:1\unsure{\text{ --- }|\to}1$.

\uline{Remarque} Dans la description combinatoire de $\underline{\mathrm{Ty}}$ et de $\operatorname{Nerf}_{\unsure{\bullet|0}}(\underline{\mathrm{Ty}})$, l'ensemble des ``signatures'' $\{\pm1\}$ joue un rôle assez \unsure{bidon}. Il peut être remplacé par n'importe quel ensemble $\Sigma$ (dit ``ensemble des signatures''), et \deleted{\unsure{les}} on construit alors une catégorie monoïdale (ou un monoïde catégorique) ayant comme ensemble d'objets le monoïde libre engendré par $\Sigma$, interprétés comme des intervalles d'entiers $[0,n]$, munis de ``signatures'' $\in\Sigma$ pour chaque intervalle de transition, \inserted{ou ce qui revient au $=$, des ``graphes'' \deleted{\unsure{pointés}}\inserted{orientés} $\Sigma$-pondérés de type combinatoire un \deleted{\unsure{dist}} segment $\Delta_{\unsure{i|1}}$,} ce qui revient à la donnée d'une suite \inserted{(ordonnée)} de \deleted{\unsure{des}} éléments $\deleted{\unsure{\Sigma|E}}\varepsilon_1,\varepsilon_2,\unsure{\text{ --- }|\ldots},\varepsilon_n\in\Sigma$, i.e. d'un ``mot'' $\varepsilon_1\varepsilon_2\ldots\varepsilon_n$, qu'on préférera écrire $\Delta_{\varepsilon_1}\Delta_{\varepsilon_2}\ldots\Delta_{\varepsilon_n}$. On peut définir de façon évidente les hom. \uline{croissants} de tels graphes pondérés orientés (sous-entendu : avec l'élément source et l'élément but dans la structure, à respecter par les morphismes), sans avoir besoin d'une interprétation catégorique. La fonctorialité de l'opération de composition des intervalles pondérés, pour cette notion de morphisme, est évidente. \unsure{Le|La}
